\section{Cel i zakres pracy}

Głównym celem pracy jest przeprowadzenie analizy porównawczej i ocena skuteczności wybranych metod mitygacji podatności dużych modeli językowych na ataki typu prompt injection.
Realizacja wymaga implementacji zróżnicowanych podejść do obrony, takich jak obrona na poziomie wejścia przy użyciu inżynierii promptów, wykorzystanie zewnętrznych modeli moderujących czy metoda oparta na inżynierii aktywacji.
W celu sprawdzenia bezpieczeństwa modeli należy również poddać je różnym wektorom ataku opartych na szkodliwych instrukcjach.
Rezultatem pracy będzie wskazanie optymalnych mechanizmów obronnych dla wybranych modeli językowych, zapewniających wysoki poziom bezpieczeństwa przy zachowaniu zadowalającej jakości ich działania.